\chapter{Maattheorie}\label{ch:b3:measure}

Hoe lang is een deelverzameling van $\R$? Het naïeve antwoord ---
ken aan elke verzameling een translatie-invariante lengte toe die
die van de intervallen voortzet --- is \emph{onmogelijk}: de
constructie van Vitali, aan het eind van dit hoofdstuk, levert een
verzameling zonder \omterm{def:b3:topology:connected}{samenhangende} lengte. De maattheorie is de
gedisciplineerde terugtocht: beperk je tot een rijke klasse
\emph{meetbare} verzamelingen, waarop een aftelbaar additieve
lengte bestaat en uniek is. De beloning is enorm --- de integraal
van Lebesgue (\cref{ch:b3:lebesgue}), de $L^p$-ruimten van de
functionaalanalyse en de hele moderne kansrekening
(\cref{ch:b3:probability}) rusten op de drie stellingen die we
hier bewijzen: het eenduidigheidslemma van Dynkin, de
uitbreidingsstelling van Carathéodory, en het bestaan van de
\omterm{def:b3:measure:lebesgueouter}{lebesguemaat}.

\section{$\sigma$-algebra's}

\begin{definition}\label{def:b3:measure:sigmaalgebra}
Een \emph{$\sigma$-algebra}\index{sigma-algebra@$\sigma$-algebra}
op een verzameling $X$ is een familie $\mathcal A$
deelverzamelingen die $\varnothing$ bevat en gesloten is onder
complementvorming en onder \emph{aftelbare} verenigingen (en dus
onder aftelbare doorsneden en verschillen, en ze bevat $X$). Het
paar $(X, \mathcal A)$ heet een \emph{meetbare ruimte}; de leden
van $\mathcal A$ heten \emph{meetbare verzamelingen}. Voor een
willekeurige familie $\mathcal E$ deelverzamelingen noteert
$\sigma(\mathcal E)$ de kleinste $\sigma$-algebra die $\mathcal E$
bevat (de doorsnede van alle zulke --- een doorsnede van
$\sigma$-algebra's is er weer een).
\end{definition}

\begin{definition}\label{def:b3:measure:borel}
De \emph{Borel-$\sigma$-algebra}\index{Borel sigma-algebra@%
Borel-$\sigma$-algebra} van een topologische ruimte is $\mathcal
B(X) = \sigma(\{\text{open verzamelingen}\})$. Op $\R$ wordt
$\mathcal B(\R)$ ook voortgebracht door de open intervallen, door
de gesloten intervallen, door de halfrechten $(-\infty, a]$ en
door de halfrechten met rationale eindpunten
(\cref{exo:b3:measure:1}) --- elke familie brengt de \omterm{def:b3:topology:topology}{open
verzamelingen} voort met aftelbare bewerkingen, want elke open
deelverzameling van $\R$ is een aftelbare vereniging van open
intervallen met rationale gegevens.
\end{definition}

\begin{definition}\label{def:b3:measure:dynkin}
Een \emph{$\pi$-systeem}\index{pi-systeem@$\pi$-systeem} is een
familie die gesloten is onder eindige doorsneden. Een
\emph{$\lambda$-systeem}\index{lambda-systeem@$\lambda$-systeem}
(dynkinklasse) is een familie $\mathcal D$ met: $X \in \mathcal
D$; uit $A, B \in \mathcal D$ met $A \subseteq B$ volgt $B
\setminus A \in \mathcal D$; en uit $A_n \uparrow A$ met $A_n \in
\mathcal D$ volgt $A \in \mathcal D$.
\end{definition}

\begin{theorem}[$\pi$--$\lambda$-lemma van Dynkin]
\label{thm:b3:measure:dynkin}
Bevat een $\lambda$-systeem $\mathcal D$ een $\pi$-systeem
$\mathcal P$, dan is $\mathcal D \supseteq
\sigma(\mathcal P)$.\index{lemma van Dynkin}
\end{theorem}

\begin{proof}
Zij $\mathcal D_0$ het kleinste $\lambda$-systeem dat $\mathcal P$
bevat (de doorsnede van alle zulke); het volstaat aan te tonen dat
$\mathcal D_0$ een $\sigma$-algebra is, want dan is
$\sigma(\mathcal P) \subseteq \mathcal D_0 \subseteq \mathcal D$.
Een $\lambda$-systeem dat gesloten is onder eindige doorsneden
\emph{is} een $\sigma$-algebra: complementen ($X \setminus A$ met
$A \subseteq X$), eindige verenigingen ($A \cup B = X \setminus
((X\setminus A)\cap(X\setminus B))$), en aftelbare verenigingen
via $\bigcup_{k \leq n}A_k \uparrow \bigcup_kA_k$. We bewijzen dus
dat $\mathcal D_0$ een $\pi$-systeem is, in twee stappen. Zij
\[
\mathcal D_1 = \{A \in \mathcal D_0 : A \cap P \in \mathcal
D_0 \ \forall P \in \mathcal P\}.
\]
$\mathcal D_1$ is een $\lambda$-systeem (alle drie de axioma's
gaat men na door met $P$ te snijden: bijvoorbeeld $(B\setminus
A)\cap P = (B \cap P)\setminus(A \cap P)$, een echt verschil
binnen $\mathcal D_0$) en bevat $\mathcal P$ ($\pi$-systeem): dus
$\mathcal D_1 = \mathcal D_0$. Zij nu
\[
\mathcal D_2 = \{A \in \mathcal D_0 : A \cap D \in \mathcal
D_0\ \forall D \in \mathcal D_0\}.
\]
Wegens de vorige stap is $\mathcal D_2 \supseteq \mathcal P$; en
$\mathcal D_2$ is een $\lambda$-systeem volgens dezelfde
verificatie: dus $\mathcal D_2 = \mathcal D_0$, wat precies zegt
dat $\mathcal D_0$ gesloten is onder doorsneden.
\end{proof}

\section{Maten}

\begin{definition}\label{def:b3:measure:measure}
Een \emph{maat}\index{maat} op $(X, \mathcal A)$ is een
afbeelding $\mu \colon \mathcal A \to [0, +\infty]$ met
$\mu(\varnothing) = 0$ die \emph{$\sigma$-additief} is: voor
paarsgewijs disjuncte $(A_n)_{n\in\N}$ geldt
\[
\mu\Bigl(\bigsqcup_n A_n\Bigr) = \sum_n \mu(A_n).
\]
$(X, \mathcal A, \mu)$ heet een \emph{maatruimte}; $\mu$ heet
\emph{eindig} als $\mu(X) < \infty$, een \emph{kansmaat} als
$\mu(X) = 1$, en \emph{$\sigma$-eindig} als $X$ een aftelbare
vereniging is van verzamelingen van eindige maat. Voorbeelden: de
telmaat op $(\N, \mathcal P(\N))$; de diracmassa $\delta_a(A) =
\mathbf 1_{a \in A}$; en, het onderwerp van dit hoofdstuk, de
\omterm{def:b3:measure:lebesgueouter}{lebesguemaat}.
\end{definition}

\begin{proposition}\label{prop:b3:measure:basics}
Zij $\mu$ een \omterm{def:b3:measure:measure}{maat}. (a) Monotonie: uit $A \subseteq B$ volgt
$\mu(A) \leq \mu(B)$. (b) Aftelbare subadditiviteit:
$\mu(\bigcup A_n) \leq \sum\mu(A_n)$. (c) \omterm{def:b3:topology:continuity}{Continuïteit} van
onderen: uit $A_n \uparrow A$ volgt $\mu(A_n) \to \mu(A)$. (d)
\omterm{def:b3:topology:continuity}{Continuïteit} van boven: uit $A_n \downarrow A$ \emph{met $\mu(A_1)
< \infty$} volgt $\mu(A_n) \to \mu(A)$.
\end{proposition}

\begin{proof}
(a) $B = A \sqcup (B\setminus A)$. (b) Maak disjunct: de $B_n =
A_n \setminus \bigcup_{k<n}A_k$ zijn disjunct met dezelfde
vereniging, en $\mu(B_n) \leq \mu(A_n)$. (c) $A = \bigsqcup_n (A_n
\setminus A_{n-1})$ (met $A_0 = \varnothing$): de partiële sommen
van $\sum\mu(A_n\setminus A_{n-1})$ zijn $\mu(A_n)$. (d) Pas (c)
toe op $A_1 \setminus A_n \uparrow A_1 \setminus A$ en trek af van
$\mu(A_1)$ --- de eindigheid maakt het aftrekken geoorloofd.
Tegenvoorbeeld zonder haar: $A_n = [n, \infty)$ voor de
\omterm{def:b3:measure:lebesgueouter}{lebesguemaat}: $A_n \downarrow \varnothing$ terwijl $\mu(A_n) =
\infty$.
\end{proof}

\begin{theorem}[Eenduidigheid]\label{thm:b3:measure:uniqueness}
Zij $\mu, \nu$ \omterm{def:b3:measure:measure}{maten} op $\sigma(\mathcal P)$, met $\mathcal P$ een
$\pi$-systeem, en $\mu = \nu$ op $\mathcal P$. Zijn er
verzamelingen $P_k \in \mathcal P$ met $P_k \uparrow X$ en
$\mu(P_k) < \infty$, dan is $\mu = \nu$ op heel
$\sigma(\mathcal P)$.
\end{theorem}

\begin{proof}
Leg $k$ vast en beschouw de eindige \omterm{def:b3:measure:measure}{maten} $\mu_k(A) = \mu(A \cap
P_k)$ en $\nu_k(A) = \nu(A \cap P_k)$ op $\sigma(\mathcal P)$: ze
vallen samen op $\mathcal P$, want $P \cap P_k \in \mathcal P$
($\pi$-systeem), en ze kennen aan $X$ dezelfde eindige waarde
$\mu(P_k)$ toe. De klasse $\mathcal D = \{A : \mu_k(A) =
\nu_k(A)\}$ is een $\lambda$-systeem: $X \in \mathcal D$; echte
verschillen door aftrekken (eindige waarden); en stijgende
limieten door de \omterm{def:b3:topology:continuity}{continuïteit} van onderen
(\cref{prop:b3:measure:basics}(c)). Ze bevat het $\pi$-systeem
$\mathcal P$, dus geeft Dynkin (\cref{thm:b3:measure:dynkin}) dat
$\mathcal D \supseteq \sigma(\mathcal P)$: $\mu_k = \nu_k$ overal.
Ten slotte geeft voor elke $A \in \sigma(\mathcal P)$ de
\omterm{def:b3:topology:continuity}{continuïteit} van onderen langs $A \cap P_k \uparrow A$ dat $\mu(A)
= \lim_k\mu_k(A) = \lim_k\nu_k(A) = \nu(A)$.
\end{proof}

\section{Buitenmaten en de stelling van Carathéodory}

\begin{definition}\label{def:b3:measure:outer}
Een \emph{buitenmaat}\index{buitenmaat} op $X$ is een afbeelding
$\mu^* \colon \mathcal P(X) \to [0, \infty]$ met
$\mu^*(\varnothing) = 0$, monotoon en aftelbaar subadditief. Een
verzameling $A$ heet \emph{$\mu^*$-meetbaar} (Carathéodory) als ze
elke verzameling additief splitst:
\[
\mu^*(E) = \mu^*(E \cap A) + \mu^*(E \setminus A)
\qquad \text{voor elke } E \subseteq X
\]
($\leq$ geldt altijd wegens de subadditiviteit; de inhoud zit in
$\geq$).
\end{definition}

\begin{theorem}[Carathéodory]\label{thm:b3:measure:caratheodory}
De $\mu^*$-meetbare verzamelingen vormen een $\sigma$-algebra
$\mathcal M$, en $\mu^*\restriction_{\mathcal M}$ is een \omterm{def:b3:measure:measure}{maat}.
Bovendien behoort elke verzameling met $\mu^*(N) = 0$ tot
$\mathcal M$ (de \omterm{def:b3:measure:measure}{maat} is \emph{\omterm{def:b3:complete:complete}{volledig}}).\index{stelling van
Carathéodory}
\end{theorem}

\begin{proof}
$\mathcal M$ bevat $\varnothing$ en is gesloten onder
complementvorming (de definiërende voorwaarde is symmetrisch in
$A$ en $X\setminus A$). \emph{Eindige verenigingen}: zij $A, B \in
\mathcal M$ en $E$ willekeurig; splits $E$ met $A$ en daarna elk
stuk met $B$:
\[
\mu^*(E) = \mu^*(E\cap A\cap B) + \mu^*(E\cap A\setminus B) +
\mu^*(E\cap B\setminus A) + \mu^*(E\setminus(A\cup B)).
\]
De eerste drie stukken overdekken $E \cap (A \cup B)$, dus geeft
de subadditiviteit $\mu^*(E) \geq \mu^*(E\cap(A\cup B)) +
\mu^*(E\setminus(A\cup B))$: dus $A \cup B \in \mathcal M$. Per
inductie volgen de eindige verenigingen; en met de complementen
zijn eindige disjunctheidsmanipulaties beschikbaar.

\emph{Additiviteit op $\mathcal M$}: voor disjuncte $A, B \in
\mathcal M$ en willekeurige $E$ is $\mu^*(E\cap(A\sqcup B)) =
\mu^*(E\cap A) + \mu^*(E \cap B)$ (splits met $A$); en per
inductie
\[
\mu^*\Bigl(E \cap \bigsqcup_{k\leq n}A_k\Bigr)
= \sum_{k\leq n}\mu^*(E\cap A_k).
\tag{$*$}
\]

\emph{Aftelbare verenigingen}: zij $(A_k) \subseteq \mathcal M$
disjunct (dat volstaat, door binnen de algebra $\mathcal M$
disjunct te maken), $A = \bigsqcup A_k$ en $E$ willekeurig. Met
$\bigsqcup_{k\leq n}A_k \in \mathcal M$ en de monotonie:
\[
\mu^*(E) = \mu^*\Bigl(E\cap\bigsqcup_{k\leq n}A_k\Bigr) +
\mu^*\Bigl(E\setminus\bigsqcup_{k\leq n}A_k\Bigr)
\geq \sum_{k \leq n}\mu^*(E\cap A_k) + \mu^*(E\setminus A)
\]
volgens ($*$). Laat $n \to \infty$ gaan en gebruik de aftelbare
subadditiviteit in de andere richting:
\[
\mu^*(E) \geq \sum_{k}\mu^*(E\cap A_k) + \mu^*(E\setminus A)
\geq \mu^*(E \cap A) + \mu^*(E\setminus A) \geq \mu^*(E):
\]
alle ongelijkheden zijn gelijkheden. Dat bewijst zowel $A \in
\mathcal M$ als, met $E = A$, de aftelbare additiviteit van
$\mu^*$ op $\mathcal M$.

\emph{Nulverzamelingen}: is $\mu^*(N) = 0$, dan is voor elke $E$:
$\mu^*(E \cap N) + \mu^*(E\setminus N) \leq 0 + \mu^*(E)$, dus $N
\in \mathcal M$.
\end{proof}

\section{De lebesguemaat op \texorpdfstring{$\R$}{R}}

\begin{definition}\label{def:b3:measure:lebesgueouter}
De \emph{lebesgue-buitenmaat}\index{lebesguemaat} van $A \subseteq
\R$ is
\[
\lambda^*(A) = \inf\Bigl\{\sum_{n} (b_n - a_n) :
A \subseteq \bigcup_n \intoo{a_n}{b_n}\Bigr\}
\]
(aftelbare overdekkingen met open intervallen).
\end{definition}

\begin{lemma}\label{lem:b3:measure:outerbasics}
$\lambda^*$ is een \omterm{def:b3:measure:outer}{buitenmaat}, invariant onder translaties, en
$\lambda^*(I) = \ell(I)$ (de lengte) voor elk interval $I$.
\end{lemma}

\begin{proof}
\omterm{def:b3:measure:outer}{Buitenmaat}: $\varnothing$ wordt door willekeurig kleine
intervallen overdekt; de monotonie is duidelijk; en de
subadditiviteit: gegeven overdekkingen van elke $A_n$ tot op
$\varepsilon 2^{-n}$ van het infimum, overdekt hun vereniging
$\bigcup A_n$ met totale lengte $\leq \sum \lambda^*(A_n) +
\varepsilon$. Translatie-invariantie: verschuif de
overdekkingen.

Lengte: het volstaat $I = \intcc ab$ te behandelen (de andere
types verschillen in de eindpunten, die \omterm{def:b3:measure:outer}{buitenmaat} $0$ hebben:
overdek met piepkleine intervallen; en vergelijk dan met
$\intcc{a+\varepsilon}{b - \varepsilon} \subseteq \intoo ab$).
$\lambda^*(\intcc ab) \leq b - a$: overdek met
$\intoo{a-\varepsilon}{b+\varepsilon}$. Zij omgekeerd $\intcc ab
\subseteq \bigcup_n\intoo{a_n}{b_n}$: wegens de
\emph{\omterm{def:b3:topology:compact}{compactheid}} (Borel--Lebesgue,
\cref{thm:b3:topology:metriccompact}) volstaan eindig veel
intervallen, zeg $I_1, \dots, I_N$. We tonen $\sum_{k\leq N}(b_k -
a_k) \geq b - a$ aan met inductie naar $N$: kies $I_{k_1} \ni a$;
is $b_{k_1} > b$, dan zijn we klaar ($b_{k_1} - a_{k_1} > b - a$);
anders wordt het segment $\intcc{b_{k_1}}b$ door de overige $N -
1$ intervallen overdekt, en geeft de inductie $\sum_{k \neq
k_1}(b_k - a_k) \geq b - b_{k_1}$, terwijl $b_{k_1} - a_{k_1} >
b_{k_1} - a$: tel op.
\end{proof}

\begin{theorem}[Lebesguemaat]\label{thm:b3:measure:lebesgue}
Elke borelverzameling van $\R$ is $\lambda^*$-meetbaar. De
beperking $\lambda$ van $\lambda^*$ tot de $\sigma$-algebra
$\mathcal L = \mathcal M_{\lambda^*} \supseteq \mathcal B(\R)$ (de
\emph{lebesgue-$\sigma$-algebra}) is de unieke \omterm{def:b3:measure:measure}{maat} op $\mathcal
B(\R)$ die aan elk interval zijn lengte toekent; ze is
translatie-invariant en $\sigma$-eindig.
\end{theorem}

\begin{proof}
Volgens \cref{thm:b3:measure:caratheodory} volstaat het aan te
tonen dat elke halfrechte $A = \intoo{-\infty}c$
$\lambda^*$-meetbaar is (de halfrechten brengen $\mathcal B$
voort, \cref{def:b3:measure:borel}). Zij $E \subseteq \R$ met
$\lambda^*(E) < \infty$ en $\bigcup I_n \supseteq E$ een
overdekking met $\sum\ell(I_n) \leq \lambda^*(E) + \varepsilon$.
Elke $I_n$ splitst in de twee intervallen $I_n' = I_n \cap A$ en
$I_n'' = I_n\setminus A$ (een interval min een halfrechte is een
interval) met $\ell(I_n') + \ell(I_n'') = \ell(I_n)$; de $I_n'$
overdekken $E \cap A$ en de $I_n''$ overdekken $E \setminus A$
(vergroot elk tot een open interval van lengte $\ell +
\varepsilon2^{-n}$ om binnen de definitie te blijven), dus
\[
\lambda^*(E\cap A) + \lambda^*(E\setminus A)
\leq \sum_n\bigl(\ell(I_n') + \ell(I_n'')\bigr) + 2\varepsilon
\leq \lambda^*(E) + 3\varepsilon .
\]
Eenduidigheid: twee \omterm{def:b3:measure:measure}{maten} die op het $\pi$-systeem van de
intervallen $\intoc ab$ met de lengte samenvallen (en er eindig
zijn) vallen samen op $\sigma(\text{intervallen}) = \mathcal B$,
volgens \cref{thm:b3:measure:uniqueness} met $P_k = \intoc{-k}k$.
$\sigma$-eindigheid: $\R = \bigcup(-k, k]$.
\end{proof}

\begin{theorem}[Regulariteit]\label{thm:b3:measure:regularity}
Voor elke $A \in \mathcal L$ geldt\index{regulariteit van de
lebesguemaat}
\[
\lambda(A) = \inf\{\lambda(U) : U \supseteq A \text{ open}\}
= \sup\{\lambda(K) : K \subseteq A \text{ compact}\}.
\]
\end{theorem}

\begin{proof}
\emph{Van buiten}: een overdekking $\bigcup I_n$ met
$\sum\ell(I_n) \leq \lambda(A) + \varepsilon$ is een \omterm{def:b3:topology:topology}{open
verzameling} $U \supseteq A$ met $\lambda(U) \leq \lambda(A) +
\varepsilon$ (subadditiviteit); is $\lambda(A) = \infty$, dan is
de uitspraak triviaal. \emph{Van binnen}: zij $A$ eerst begrensd,
$A \subseteq [-M, M]$. Kies een open $U \supseteq ([-M,M]\setminus
A)$ met $\lambda(U) \leq \lambda([-M,M]\setminus A) +
\varepsilon$; dan is $K = [-M, M]\setminus U$ \omterm{def:b3:topology:compact}{compact} met $K
\subseteq A$, en
\[
\lambda(K) = \lambda([-M,M]) - \lambda([-M,M]\cap U)
\geq \lambda([-M,M]) - \bigl(\lambda([-M,M]) - \lambda(A) +
\varepsilon\bigr) = \lambda(A) - \varepsilon .
\]
Voor algemene $A$: $\lambda(A) = \lim_M\lambda(A\cap[-M,M])$
(\omterm{def:b3:topology:continuity}{continuïteit} van onderen), en pas daarbinnen het begrensde geval
toe.
\end{proof}

\begin{example}\label{ex:b3:measure:cantor}
De cantorverzameling (\cref{exo:b3:topology:10}) heeft
$\lambda(C) = 0$: $C \subseteq C_n$, een vereniging van $2^n$
intervallen van lengte $3^{-n}$, dus $\lambda(C) \leq (2/3)^n \to
0$. Een overaftelbare nulverzameling --- de kardinaliteit ziet de
\omterm{def:b3:measure:measure}{maat} niet. Omgekeerd zijn \emph{dikke cantorverzamelingen}
(\cref{exo:b3:measure:5}) nergens \omterm{def:b3:topology:interior}{dicht} met positieve \omterm{def:b3:measure:measure}{maat}: ook
de \omterm{def:b3:topology:topology}{topologie} ziet de \omterm{def:b3:measure:measure}{maat} niet. De weekendopgave drijft dat
samenspel tot zijn treffende conclusie: er bestaan
lebesgue-meetbare verzamelingen die geen borelverzameling zijn.
\end{example}

\begin{theorem}[Vitali]\label{thm:b3:measure:vitali}
Er bestaat geen \omterm{def:b3:measure:measure}{maat} op \emph{alle} deelverzamelingen van $\R$ die
translatie-invariant is en aan elk interval zijn lengte toekent.
In het bijzonder is $\mathcal L \neq \mathcal P(\R)$: er bestaan
niet-meetbare verzamelingen.\index{niet-meetbare verzameling van
Vitali}
\end{theorem}

\begin{proof}
Stel dat $\mu$ er een was. Beschouw op $\intcc01$ de equivalentie
$x \sim y \iff x - y \in \Q$; kies met het \emph{keuzeaxioma} per
klasse één representant in $\intcc01$: een verzameling $V$. Voor
$q \in \Q\cap\intcc{-1}1$ zijn de translaten $V + q$ paarsgewijs
disjunct (twee punten van $V$ die een rationaal getal verschillen
zouden equivalent en toch verschillende representanten zijn), en
\[
\intcc01 \subseteq \bigsqcup_{q \in \Q\cap\intcc{-1}1}(V + q)
\subseteq \intcc{-1}2 :
\]
de eerste inclusie omdat elke $x \in \intcc01$ van haar
representant $v$ verschilt met een rationaal getal $q = x - v \in
\intcc{-1}1$. De monotonie en de $\sigma$-additiviteit geven
\[
1 \leq \sum_{q}\mu(V + q) \leq 3,
\qquad\text{met } \mu(V + q) = \mu(V) \text{ voor alle } q .
\]
Een oneindige som van de constante $\mu(V)$ is $0$ of $\infty$:
beide grenzen kunnen niet tegelijk gelden. Zo'n $\mu$ bestaat dus
niet --- en $V \notin \mathcal L$, want $\lambda$ heeft op
$\mathcal L$ alle gebruikte eigenschappen.
\end{proof}

\begin{remark}\label{rem:b3:measure:banachtarski}
In $\R^3$ is het falen dramatischer: de paradox van
Banach--Tarski ontbindt een bol in vijf stukken die zich, door
rotaties en translaties, weer laten samenvoegen tot \emph{twee}
bollen van dezelfde straal --- zodat er zelfs geen eindig
additief, rotatie-invariant volume op alle deelverzamelingen van
$\R^3$ bestaat. De stukken zijn uiteraard niet meetbaar.
Meetbaarheid is geen bureaucratische voorzichtigheid; het is de
grens van de \omterm{def:b3:topology:connected}{samenhang}.
\end{remark}

\begin{method}\label{met:b3:measure:goodsets}
Het \emph{principe van de goede verzamelingen}: om te bewijzen dat
alle verzamelingen van $\sigma(\mathcal E)$ een eigenschap hebben,
toon je aan dat de goede verzamelingen een $\sigma$-algebra vormen
(of een $\lambda$-systeem, als de eigenschap maattheoretisch is en
$\mathcal E$ een $\pi$-systeem --- dan Dynkin) die $\mathcal E$
bevat. Bijna elk bewijs in dit hoofdstuk en het volgende is er een
geval van. Om twee \omterm{def:b3:measure:measure}{maten} gelijk te bewijzen: controleer ze op een
voortbrengend $\pi$-systeem, plus $\sigma$-eindigheid
(\cref{thm:b3:measure:uniqueness}). Om een \omterm{def:b3:measure:measure}{maat} te bouwen: bouw
met overdekkingen een \omterm{def:b3:measure:outer}{buitenmaat} en roep Carathéodory in.
\end{method}

\section{Oefeningen}

\begin{exercise}[$\star$]\label{exo:b3:measure:1}
(a) Toon aan dat $\{A \subseteq X : A$ of $X\setminus A$ is
aftelbaar$\}$ een $\sigma$-algebra is: die voortgebracht door de
singletons.
(b) Toon aan dat $\mathcal B(\R)$ wordt voortgebracht door elk
van: de open intervallen; de gesloten intervallen; de halfrechten
$\intoc{-\infty}a$; en de halfrechten met $a \in \Q$.
(c) Is de familie van de \emph{eindige disjuncte verenigingen van
intervallen} $\intoc ab$ een $\sigma$-algebra? Een algebra
(gesloten onder complementvorming en eindige verenigingen)?
\end{exercise}

\begin{exercise}[$\star$]\label{exo:b3:measure:2}
(a) Bewijs de in-en-uitsluiting voor een eindige \omterm{def:b3:measure:measure}{maat}: $\mu(A\cup
B) = \mu(A) + \mu(B) - \mu(A\cap B)$, en de versie met drie
verzamelingen.
(b) Geef een voorbeeld dat laat zien dat de \omterm{def:b3:topology:continuity}{continuïteit} van boven
(\cref{prop:b3:measure:basics}(d)) zonder de
eindigheidsveronderstelling faalt.
(c) Toon aan dat een aftelbare verzameling \omterm{def:b3:measure:lebesgueouter}{lebesguemaat} nul heeft.
Leid af dat $\lambda(\Q) = 0$ en $\lambda(\intcc01\setminus\Q) =
1$.
\end{exercise}

\begin{exercise}[$\star\star$]\label{exo:b3:measure:3}
Zij $\mu, \nu$ kansmaten op $\mathcal B(\R)$ met
$\mu(\intoc{-\infty}t) = \nu(\intoc{-\infty}t)$ voor alle $t \in
\R$. Toon aan dat $\mu = \nu$. (Daarmee is de
\emph{verdelingsfunctie} $F(t) = \mu(\intoc{-\infty}t)$ een
volledige invariant --- de grondslag van
\cref{ch:b3:probability}.)
\end{exercise}

\begin{exercise}[$\star\star$]\label{exo:b3:measure:4}
(Borel--Cantelli, maatversie) Zij $(A_n)$ meetbaar met
$\sum_n\mu(A_n) < \infty$, en $\limsup A_n =
\bigcap_N\bigcup_{n\geq N}A_n$ (de punten die tot oneindig veel
$A_n$ behoren). Toon aan dat $\mu(\limsup A_n) = 0$. Toepassing:
voor bijna elke $x \in \intcc01$ voldoen slechts eindig veel $n$
aan $\abs{x - p/q_n} \leq 4^{-n}$ voor het $n$-de rationale getal
$p/q_n$ van een opsomming van $\Q\cap\intcc01$.
\end{exercise}

\begin{exercise}[$\star\star$]\label{exo:b3:measure:5}
(Dikke cantorverzameling) Herhaal de cantorconstructie op
$\intcc01$, maar haal in stap $n$ uit elk van de $2^{n-1}$
intervallen slechts een \emph{gecentreerd} open interval van
lengte $4^{-n}$ weg. Toon aan dat de resulterende $K = \bigcap
K_n$ \omterm{def:b3:topology:compact}{compact} is, een leeg inwendige heeft (geen enkel interval
overleeft) en
\[
\lambda(K) = 1 - \sum_{n\geq1}2^{n-1}4^{-n} = \tfrac12 :
\]
een nergens \omterm{def:b3:topology:interior}{dichte verzameling} van \omterm{def:b3:measure:measure}{maat} $\frac12$. Leid daaruit
een \emph{\omterm{rem:b3:complete:meagre}{magere}} deelverzameling van $\intcc01$ van volle \omterm{def:b3:measure:measure}{maat}
$1$ af, en een open \omterm{def:b3:topology:interior}{dichte} deelverzameling van \omterm{def:b3:measure:measure}{maat} $<
\varepsilon$.
\end{exercise}

\begin{exercise}[$\star\star$]\label{exo:b3:measure:6}
Zij $\mu$ een \omterm{def:b3:measure:measure}{maat} op $\mathcal B(\R)$, invariant onder
translaties, met $c = \mu(\intoc01) < \infty$. Toon aan dat $\mu =
c\,\lambda$ op $\mathcal B(\R)$. \emph{(Bereken $\mu$ op dyadische
intervallen door $\intoc01$ in $2^n$ translaten te verdelen, en
roep dan \cref{thm:b3:measure:uniqueness} in.)}
\end{exercise}

\begin{exercise}[$\star\star$]\label{exo:b3:measure:7}
(Benadering) Zij $A \in \mathcal L$ met $\lambda(A) < \infty$ en
$\varepsilon > 0$. Toon aan dat er een \emph{eindige} vereniging
van intervallen $B$ is met $\lambda(A\,\triangle\,B) <
\varepsilon$ ($\triangle$ is het symmetrische verschil).
\emph{(Regulariteit: klem $K \subseteq A \subseteq U$ in en
gebruik de structuur van de open $U$ als aftelbare vereniging van
intervallen, plus de \omterm{def:b3:topology:compact}{compactheid} van $K$.)}
\end{exercise}

\begin{exercise}[$\star\star\star$]\label{exo:b3:measure:8}
(Steinhaus) Zij $A \in \mathcal L$ met $\lambda(A) > 0$. Toon aan
dat $A - A = \{x - y : x, y \in A\}$ een interval rond $0$ bevat.
\emph{(Herleid tot $\lambda(A) < \infty$; vind met regulariteit in
de trant van \cref{exo:b3:measure:7} een interval $I$ met
$\lambda(A \cap I) > \frac34\ell(I)$; dan liggen voor $\abs t <
\frac12\ell(I)$ de verzamelingen $A\cap I$ en $(A\cap I) + t$
beide in een interval van lengte $\frac32\ell(I)$ met totale \omterm{def:b3:measure:measure}{maat}
$> \frac32\ell(I)$: ze moeten elkaar snijden.)}
\end{exercise}

\begin{exercise}[$\star\star\star$]\label{exo:b3:measure:9}
Toon aan dat elke $A \in \mathcal L$ met $\lambda(A) > 0$ een
niet-meetbare deelverzameling bevat. \emph{(Snijd $A$ met de
translaten $V + q$ van de verzameling van Vitali: waren alle $A
\cap (V+q)$ meetbaar, dan was elk van hen een nulverzameling
volgens het argument van \cref{thm:b3:measure:vitali} ---
Steinhaus (\cref{exo:b3:measure:8}) helpt: een meetbare
verzameling van positieve \omterm{def:b3:measure:measure}{maat} binnen $V + q$ zou geven dat
$(V+q) - (V+q)$ een interval bevat, in strijd met het feit dat die
verschilverzameling $\Q$ alleen in $0$ snijdt; besluit met de
subadditiviteit.)}
\end{exercise}

\begin{exercise}[$\star\star$]\label{exo:b3:measure:10}
Toon aan dat een $A \subseteq \R$ met $\lambda^*(A) < \infty$
lebesgue-meetbaar is dan en slechts dan als er voor elke
$\varepsilon > 0$ een open $U \supseteq A$ is met $\lambda^*(U
\setminus A) < \varepsilon$, en dan en slechts dan als er een
$G_\delta$-verzameling $G \supseteq A$ is met $\lambda^*(G\setminus
A) = 0$. (Lebesgueverzamelingen zijn dus borelverzamelingen op
nulverzamelingen na.)
\end{exercise}

\begin{exercise}[$\star\star$]\label{exo:b3:measure:11}
(\omterm{def:b3:topology:continuity}{Continuïteit} langs monotone limieten, en de scherpte ervan)
(a) Toon aan dat voor meetbare verzamelingen $\mu(\liminf A_n)
\leq \liminf\mu(A_n)$ (Fatou voor verzamelingen), en dat, als
$\mu\bigl(\bigcup A_n\bigr) < \infty$, ook $\limsup\mu(A_n) \leq
\mu(\limsup A_n)$.
(b) Geef, voor de \omterm{def:b3:measure:lebesgueouter}{lebesguemaat} op $\R$, een rij met $\mu(A_n) = 1$
voor alle $n$ en toch $\mu(\limsup A_n) = 0$: de
eindigheidshypothese in de tweede ongelijkheid is geen
versiering.
(c) Leid af: uit $\sum\mu(A_n) < \infty$ volgt $\mu(\limsup A_n) =
0$ (opnieuw Borel--Cantelli), en als de $A_n$ stijgen of dalen
(met $\mu(A_1) < \infty$ in het dalende geval), dan is $\mu(\lim
A_n) = \lim\mu(A_n)$.
\end{exercise}

\begin{exercise}[$\star\star\star$]\label{exo:b3:measure:12}
(Stelling van Egorov) Zij $\mu(X) < \infty$ en $f_n \to f$
puntsgewijs, alle meetbaar (reëelwaardig). Stel voor $k, N \geq 1$
\[
E_{k,N} = \bigcap_{n \geq N}\Bigl\{x : \abs{f_n(x) - f(x)}
\leq \tfrac1k\Bigr\} .
\]
(a) Toon aan dat voor vaste $k$ geldt $E_{k,N} \nearrow X$ als $N
\to \infty$, en leid een $N_k$ af met $\mu(X \setminus E_{k,N_k})
\leq \varepsilon2^{-k}$.
(b) Besluit de \emph{stelling van Egorov}: voor elke $\varepsilon
> 0$ is er een meetbare $A$ met $\mu(X\setminus A) \leq
\varepsilon$ zodanig dat $f_n \to f$ \emph{uniform op $A$} ---
puntsgewijze convergentie is uniforme convergentie buiten een
willekeurig kleine verzameling.
(c) Toon aan dat de stelling faalt op $(\R, \lambda)$: de
bewegende bulten $f_n = \mathbf 1_{\intcc n{n+1}}$ convergeren
puntsgewijs naar $0$ maar op geen enkel complement van een
verzameling van eindige \omterm{def:b3:measure:measure}{maat} uniform. Waar gebruikte (a) dat
$\mu(X) < \infty$?
\end{exercise}

\section{Probleem: de cantortrap en een meetbare verzameling die
geen borelverzameling is}

\begin{omfigure}
\begin{tikzpicture}[scale=5.2]
  \draw[->] (0,0) -- (1.06,0) node[right, font=\small] {$x$};
  \draw[->] (0,0) -- (0,1.06);
  \draw[black!30] (1,0) -- (1,1) -- (0,1);
  \draw[omThm, thick]
    (0,0)
    -- (0.037,0.125) -- (0.074,0.125)
    -- (0.111,0.25)  -- (0.222,0.25)
    -- (0.259,0.375) -- (0.296,0.375)
    -- (0.333,0.5)   -- (0.667,0.5)
    -- (0.704,0.625) -- (0.741,0.625)
    -- (0.778,0.75)  -- (0.889,0.75)
    -- (0.926,0.875) -- (0.963,0.875)
    -- (1,1);
  \node[font=\small, omThm] at (0.32,0.86)
    {$c$ (benadering op diepte 3)};
\end{tikzpicture}

{\small De cantortrap van Cantor en Vitali: constant op elk gat
van de cantorverzameling, en toch \omterm{def:b3:topology:continuity}{continu} klimmend van $0$ naar
$1$. Haar afgeleide verdwijnt bijna overal --- al het klimmen
gebeurt op een nulverzameling.}
\end{omfigure}

\begin{problem}[{Weekendopgave --- de duivelstrap, en $\mathcal
B(\R) \subsetneq \mathcal L$}]
\label{pb:b3:measure:1}
We construeren de \emph{functie van Cantor en Vitali} (de
duivelstrap), gebruiken haar om \omterm{def:b3:measure:measure}{maat} op pathologische wijze te
vervoeren, en besluiten met een stelling die geen zacht argument
oplevert: er bestaan lebesgue-meetbare verzamelingen die geen
borelverzameling zijn. Notatie: $C$ is de cantorverzameling,
$C_n$ haar $n$-de stadium ($2^n$ intervallen van lengte $3^{-n}$),
en elke $x \in C$ heeft ternaire cijfers $x = \sum 2b_n3^{-n}$ met
$b_n \in \{0,1\}$ (\cref{exo:b3:topology:10}).

\medskip
\noindent\textbf{Deel I --- De trap.} Definieer $c_0(x) = x$ en
$c_{n+1}$ uit $c_n$ door
\[
c_{n+1}(x) = \begin{cases}
\tfrac12\,c_n(3x) & 0 \leq x \leq \tfrac13,\\[2pt]
\tfrac12 & \tfrac13 \leq x \leq \tfrac23,\\[2pt]
\tfrac12 + \tfrac12\,c_n(3x - 2) & \tfrac23 \leq x \leq 1.
\end{cases}
\]
\begin{enumerate}
  \item Toon aan dat elke $c_n$ \omterm{def:b3:topology:continuity}{continu} en niet-dalend is met
        $c_n(0) = 0$ en $c_n(1) = 1$, en dat $\norm{c_{n+1} -
        c_n}_\infty \leq \tfrac12\norm{c_n - c_{n-1}}_\infty$.
  \item Leid af dat $(c_n)$ uniform convergeert naar een \omterm{def:b3:topology:continuity}{continue},
        niet-dalende $c$ met $c(0) = 0$ en $c(1) = 1$ (\emph{de}
        functie van Cantor en Vitali), die aan dezelfde
        zelfgelijkvormige betrekkingen voldoet als de $c_{n+1}$
        hierboven.
  \item Toon aan dat $c$ constant is op elke \omterm{def:b3:topology:components}{samenhangscomponent}
        van $\intcc01\setminus C$, en dat voor $x = \sum_n
        2b_n3^{-n} \in C$ geldt $c(x) = \sum_n b_n2^{-n}$: de trap
        leest cantorcijfers binair (de functie $g$ van
        \cref{pb:b3:topology:1}, monotoon en globaal gemaakt).
  \item Leid af dat $c$ in elk punt van $\intcc01\setminus C$
        differentieerbaar is met $c' = 0$: dus $c' = 0$
        \emph{$\lambda$-bijna overal}
        (\cref{ex:b3:measure:cantor}). Besluit dat de
        hoofdstelling van de integraalrekening voor $c$ faalt:
        \[
        c(1) - c(0) = 1 \neq 0 = \int_0^1 c'(t)\,\dd t
        \]
        (de integraal over de verzameling van volle \omterm{def:b3:measure:measure}{maat} waar $c'
        = 0$; vooruitlopend op \cref{ch:b3:lebesgue}:
        nulverzamelingen beïnvloeden integralen niet). Welke
        hypothese van de $\mathcal C^1$-hoofdstelling wordt
        geschonden?
  \item Toon aan dat $c(C) = \intcc01$: de nulverzameling $C$
        wordt \emph{op} een verzameling van volle \omterm{def:b3:measure:measure}{maat} afgebeeld.
\end{enumerate}

\medskip
\noindent\textbf{Deel II --- Het scheve \omterm{def:b3:topology:continuity}{homeomorfisme}.} Zij $h(x)
= \frac{c(x) + x}{2}$.
\begin{enumerate}[resume]
  \item Toon aan dat $h \colon \intcc01 \to \intcc01$ een
        \omterm{def:b3:topology:continuity}{homeomorfisme} is (strikt stijgend, \omterm{def:b3:topology:continuity}{continu}, surjectief).
  \item Toon aan dat $\lambda\bigl(h(\intcc01\setminus C)\bigr) =
        \tfrac12$: op elk gat van lengte $\ell$ werkt $h$ als een
        affiene afbeelding met helling $\tfrac12$, en de gaten
        hebben samen lengte $1$.
  \item Leid af dat $\lambda\bigl(h(C)\bigr) = \tfrac12$: het
        \omterm{def:b3:topology:continuity}{homeomorfe} beeld van een nulverzameling kan positieve \omterm{def:b3:measure:measure}{maat}
        hebben. (Waar botst dat met de naïeve intuïtie over
        ``grootte''?)
\end{enumerate}

\medskip
\noindent\textbf{Deel III --- Een meetbare verzameling die geen
borelverzameling is.}
\begin{enumerate}[resume]
  \item Kies met \cref{exo:b3:measure:9} een niet-meetbare $W
        \subseteq h(C)$. Toon aan dat $Z = h^{-1}(W) \subseteq C$
        lebesgue-meetbaar is. \emph{(Het is een deelverzameling
        van een nulverzameling; \omterm{def:b3:complete:complete}{volledigheid},
        \cref{thm:b3:measure:caratheodory}.)}
  \item Toon aan dat het origineel van een borelverzameling onder
        een \omterm{def:b3:topology:continuity}{continue afbeelding} een borelverzameling is.
        \emph{(Principe van de goede verzamelingen: $\{B :
        h^{-1}(B) \in \mathcal B\}$ is een $\sigma$-algebra die de
        \omterm{def:b3:topology:topology}{open verzamelingen} bevat --- let op de richting van de
        afbeelding.)}
  \item Besluit dat $Z$ \emph{geen} borelverzameling is: was ze
        het, dan was $W = (h^{-1})^{-1}(Z)$ er een (pas vraag 10
        toe op de \omterm{def:b3:topology:continuity}{continue} $h^{-1}$) en dus meetbaar ---
        tegenspraak. Bijgevolg is
        \[
        \boxed{\ \mathcal B(\R) \subsetneq \mathcal L\ }
        \]
        en vergroot de \omterm{def:b3:complete:complete}{volledigheid} de borelwereld werkelijk.
  \item Geef een lebesgue-meetbare functie $g$ en een \omterm{def:b3:topology:continuity}{continue}
        functie $\varphi$ zodanig dat $g \circ \varphi$ niet
        lebesgue-meetbaar is: meetbaarheid stelt zich, anders dan
        \omterm{def:b3:topology:continuity}{continuïteit}, niet samen. \emph{(Neem $g = \mathbf 1_Z$ en
        $\varphi = h^{-1}$, vooruitlopend op de definitie van
        meetbare functies uit \cref{ch:b3:lebesgue}: originelen
        van borelverzamelingen zijn lebesgueverzamelingen. Waar
        moet je oppassen met de vraag welke $\sigma$-algebra op
        het doel wordt gebruikt?)}
\end{enumerate}

\medskip
\noindent\textbf{Deel IV --- Slotwoord.}
\begin{enumerate}[resume]
  \item Orden de volgende klassen naar strikte inclusie en
        verantwoord elke striktheid met een voorbeeld uit dit
        hoofdstuk en zijn probleem: aftelbare verzamelingen;
        borelnulverzamelingen; lebesguenulverzamelingen;
        borelverzamelingen; lebesgueverzamelingen; willekeurige
        verzamelingen.
\end{enumerate}

\medskip
\noindent\textbf{Deel V --- De cantormaat: massa op een
nulverzameling.} De trap is de verdelingsfunctie van een
opmerkelijke \omterm{def:b3:measure:measure}{maat}, die we nu met het gereedschap van dit hoofdstuk
bouwen.
\begin{enumerate}[resume]
  \item (Lebesgue--Stieltjes, bestaan) Zij $F\colon\R\to\R$
        niet-dalend, \omterm{def:b3:topology:continuity}{continu} en begrensd. Definieer op halfopen
        intervallen $\rho\bigl(\intoc ab\bigr) = F(b) - F(a)$ en,
        voor $A \subseteq \R$,
        \[
        \mu_F^*(A) = \inf\Bigl\{\sum_k\bigl(F(b_k) -
        F(a_k)\bigr) : A \subseteq
        \bigcup_k\intoc{a_k}{b_k}\Bigr\} .
        \]
        Toon aan dat $\mu_F^*$ een \omterm{def:b3:measure:outer}{buitenmaat} is en dat
        $\mu_F^*\bigl(\intoc ab\bigr) = F(b) - F(a)$ \emph{(boots
        het compactheidsargument van
        \cref{thm:b3:measure:lebesgue} na, waarbij je elke
        $\intoc{a_k}{b_k}$ tot een open interval vergroot tegen
        een $F$-kost $\leq \varepsilon2^{-k}$ --- waar wordt de
        \omterm{def:b3:topology:continuity}{continuïteit} van $F$ gebruikt?)}.
  \item Toon aan dat elke borelverzameling $\mu_F^*$-meetbaar is
        in de zin van Carathéodory \emph{(net als bij Lebesgue
        volstaat het de halfrechten te toetsen; volg het bewijs
        van de toepassing van
        \cref{thm:b3:measure:caratheodory})}, zodat $\mu_F =
        \mu_F^*$ beperkt tot $\mathcal B(\R)$ een \omterm{def:b3:measure:measure}{maat} is met
        $\mu_F(\intoc ab) = F(b) - F(a)$: de
        \emph{Lebesgue--Stieltjes-maat} van $F$.
  \item Pas dit toe op de trap ($F = c$, voortgezet met $0$ op
        $\R_-$ en $1$ op $\intco1\infty$): de \emph{cantormaat}
        $\mu$. Toon aan dat $\mu(\R) = 1$, dat elk gat van de
        cantorverzameling $\mu$-nul is ($c$ is daar constant), en
        besluit
        \[
        \mu(C) = 1, \qquad \lambda(C) = 0 :
        \]
        $\mu$ en $\lambda$ leven op disjuncte dragers ($C$ en haar
        complement). Twee \omterm{def:b3:measure:measure}{maten} in die positie heten
        \emph{onderling singulier}, genoteerd $\mu \perp
        \lambda$.
  \item Toon aan dat $\mu$ geen atomen heeft: $\mu(\{x\}) = 0$
        voor elke $x$ \emph{(de \omterm{def:b3:topology:continuity}{continuïteit} van $c$)}. Een
        atoomloze kansmaat gedragen door een lebesgue-nulle
        \omterm{def:b3:topology:compact}{compacte} verzameling: vergelijk met de tot nu toe geziene
        \omterm{def:b3:measure:measure}{maten}.
  \item (Munten gooien in vermomming) Zij voor een woord
        $(\varepsilon_1, \dots, \varepsilon_m) \in \{0,1\}^m$ de
        verzameling $C_{\varepsilon}$ die van de $x \in C$ waarvan
        de ternaire cijfers voldoen aan $b_i(x) = \varepsilon_i$
        voor $i \leq m$ (een van de $2^m$ cantorstukken van diepte
        $m$). Toon aan dat $\mu(C_\varepsilon) = 2^{-m}$ \emph{(de
        trap klimt $2^{-m}$ over dat stuk: gebruik Deel I, vraag
        3)}. De cantormaat is de wet van een oneindige rij eerlijke
        muntworpen, ternair gelezen --- \cref{ch:b3:probability}
        zal dat precies maken.
  \item Bewijs de zelfgelijkvormigheid: voor elke borelverzameling
        $A$ is
        \[
        \mu(A) = \tfrac12\,\mu(3A) +
        \tfrac12\,\mu(3A - 2) ,
        \]
        waarbij $3A - 2 = \{3x - 2 : x \in A\}$ \emph{(controleer
        het op de voortbrengende intervallen $\intoc ab$ via de
        zelfgelijkvormige betrekkingen van $c$, en roep dan de
        eenduidigheid in, \cref{thm:b3:measure:uniqueness})}.
  \item Toon aan dat de spiegeling $s(x) = 1 - x$ de \omterm{def:b3:measure:measure}{maat} $\mu$
        bewaart: $\mu(s(A)) = \mu(A)$ \emph{(via $c(1 - x) = 1 -
        c(x)$, wat uit de symmetrie van de constructie volgt ---
        bewijs dat)}.
  \item Bereken de eerste twee momenten van $\mu$, dat wil zeggen
        van een toevallig punt $X$ met wet $\mu$ (de integralen
        mogen als limieten van sommen over de stukken van diepte
        $m$ behandeld worden, vooruitlopend op
        \cref{ch:b3:lebesgue}): de symmetrie geeft $\int x\,\dd\mu
        = \frac12$, en de zelfgelijkvormigheid geeft
        \[
        \int x^2\,\dd\mu = \frac38,
        \qquad\text{en dus}\qquad
        \operatorname{Var}(X) = \frac18 .
        \]
        Vergelijk met de uniforme wet op $\intcc01$ (variantie
        $\frac1{12}$): de cantormassa, naar de randen geduwd,
        spreidt \emph{meer}.
  \item Toon aan dat de (topologische) \emph{drager} van $\mu$ ---
        de kleinste gesloten verzameling van volle \omterm{def:b3:measure:measure}{maat} --- precies
        $C$ is.
  \item (Synthese) De trap $c$ is \omterm{def:b3:topology:continuity}{continu} en niet-dalend en faalt
        toch voor de hoofdstelling van de integraalrekening (Deel
        I); en de \omterm{def:b3:measure:measure}{maat} $\mu_c$ is een kansmaat, atoomloos,
        singulier ten opzichte van $\lambda$. Leg in een korte
        alinea uit hoe dat twee gezichten van één verschijnsel
        zijn, en formuleer de algemene moraal: niet-dalende
        functies corresponderen met \omterm{def:b3:measure:measure}{maten} ($F \leftrightarrow
        \mu_F$), differentieerbaarheid bijna overal correspondeert
        met het ``absoluut \omterm{def:b3:topology:continuity}{continue} deel'', en $c$ is de
        standaardgetuige dat een \omterm{def:b3:topology:continuity}{continue} $F$ helemaal \emph{geen}
        absoluut \omterm{def:b3:topology:continuity}{continu} deel hoeft te dragen.
  \item (De scherpe continuïteitsmodulus) Zij $s = \frac{\ln
        2}{\ln 3}$. Toon aan dat $c$ hölder \omterm{def:b3:topology:continuity}{continu} is met
        exponent $s$:
        \[
        \abs{c(x) - c(y)} \leq 4\,\abs{x - y}^{s}
        \qquad (x, y \in \intcc01),
        \]
        en dat geen enkele exponent $t > s$ kan werken, ook niet
        lokaal. Leid de maattheoretische vorm af: voor elke $x$ en
        elke $r \in \intoc01$ is
        \[
        \mu\bigl(\intcc{x - r}{x + r}\bigr) \leq 8\,r^{s} .
        \]
        \emph{(Vergelijk een triadisch rooster van diepte $m$ met
        de schaal van $\abs{x - y}$; vraag 18 geeft de klim over
        elk stuk. De exponent $s$ is de hausdorffdimensie van $C$,
        zoals latere cursussen zullen zeggen.)}
  \item (Zelfgelijkvormigheid kenmerkt $\mu$) Bewijs de omkering
        van vraag 19: is $\nu$ een kansmaat op $\mathcal B(\R)$
        gedragen door $\intcc01$ die voldoet aan
        \[
        \nu(A) = \tfrac12\,\nu(3A) + \tfrac12\,\nu(3A - 2)
        \qquad (A \in \mathcal B(\R)),
        \]
        dan is $\nu = \mu$. \emph{(Herhaal de betrekking $m$ maal
        om $\nu$ over de $2^m$ cantorstukken van diepte $m$ te
        spreiden, schat $\nu(\intoc ab)$ af tegen het aantal
        stukken binnen $\intoc ab$, en laat $m \to \infty$ gaan;
        maak het af met \cref{thm:b3:measure:uniqueness}.)}
\end{enumerate}
\end{problem}
